% Smoke test: apchem-exam.cls -- MCQ engine, FRQ blocks, rubric, point totals.
\documentclass{apchem-exam}

\apcunit{0}
\apcunittitle{Infrastructure Smoke Test}
\apcdoctitle{Unit Exam}

\begin{document}
\apctitleblock

\examsection{I}{Multiple Choice}{2 questions \textbullet\ 3 minutes}{\nocalculator}

\begin{mcq}
  A sample of neon gives mass-spectrum peaks at 20, 21, and 22 u with relative
  areas 90.5, 0.3, and 9.2. The average atomic mass is closest to
  \choices
    \choice{20.0 u}
    \correct{20.2 u}
    \choice{21.0 u}
    \choice{22.0 u}
\end{mcq}
\why{Weighted average is dominated by the 20 u peak but pulled slightly up by 22 u.}
\distractor{A}{took the most abundant isotope's mass as the answer}
\distractor{C}{took the unweighted mean of the three masses}

\begin{mcq}
  Which species has the ground-state configuration $1s^2 2s^2 2p^6$?
  \choices
    \choice{\ce{Na}}
    \choice{\ce{F}}
    \correct{\ce{Mg^2+}}
    \choice{\ce{S^2-}}
\end{mcq}
\why{\ce{Mg^2+} has 10 electrons, isoelectronic with Ne.}

\examsection{II}{Free Response}{1 question \textbullet\ 9 minutes}{\calculatorok}

\frq{short}{4}{CED 1.2 Mass Spectra \textbullet\ SP 4}

A pure element gives two mass-spectrum peaks: 62.93 u (69.2\%) and 64.93 u
(30.8\%).

\begin{parts}
  \item Calculate the average atomic mass. \workspace[2.5cm]
  \item Identify the element. \answerblank[3cm]{copper}
  \item A student claims the sample contains two different elements.
        Refute this claim in one sentence. \answerlines[2]{Both peaks arise from
        atoms with the same number of protons; they are isotopes of one element.}
\end{parts}

\begin{rubric}
  \rpoint{1}{correct weighted-average setup with both abundances}
  \rpoint{1}{answer $63.55\pm0.02$ u with correct units}
  \rpoint{1}{identifies copper from the periodic table}
  \rpoint{1}{refutation references equal proton number / isotopes}
\end{rubric}

\examtotal

\end{document}
